\chapter{Gram Stain}

\section*{What Is This Test?}
The Gram stain is the most widely used differential staining procedure in clinical microbiology, providing rapid presumptive information about the presence and category of bacteria and inflammatory cells in clinical specimens \cite{forbes2019bailey}. The stain differentiates bacteria into two broad groups based on cell wall structural differences: Gram-positive organisms (thick peptidoglycan layer retaining the crystal violet-iodine complex, appearing purple/blue) and Gram-negative organisms (thin peptidoglycan layer between the inner and outer membranes; the crystal violet-iodine complex is extracted by the organic solvent decolorizer; counterstained with safranin or basic fuchsin, appearing red/pink). The Gram stain takes approximately 3--5 minutes to perform from a prepared smear and provides immediately actionable information that guides empiric antibiotic therapy hours to days before culture results are available. In some situations (e.g., CSF from suspected bacterial meningitis, joint fluid from suspected septic arthritis, urethral discharge in symptomatic males), the Gram stain can be diagnostic on its own. In respiratory specimens, Gram stain assessment of specimen adequacy (squamous epithelial cells vs. WBCs) determines whether the specimen is acceptable for culture.

\section*{Alternative Names}
\begin{itemize}
  \item Gram's Stain
  \item Bacterial Smear
  \item Direct Smear
  \item Gram Stain with Culture
  \item Gram Stain for Bacteria
  \item WBC and Squamous Epithelial Cell Assessment
\end{itemize}
\section*{Principle}
The four-step staining procedure: (1) Primary stain: The heat-fixed or methanol-fixed smear is flooded with crystal violet (a basic purple dye) for 30--60 seconds. Both gram-positive and gram-negative bacteria take up the dye. (2) Mordant: The smear is flooded with Gram's iodine (iodine-potassium iodide) for 30--60 seconds. The iodine forms a large crystal violet-iodine (CV-I) complex within the cytoplasm. (3) Decolorization: The smear is flooded with an organic solvent (95\% ethanol, acetone, or a 1:1 mixture of acetone and ethanol) for 5--10 seconds until the runoff is clear. This is the critical differential step. The thick peptidoglycan layer of gram-positive organisms (20--80 nm, 40--60 layers) resists decolorization and retains the CV-I complex. The thin peptidoglycan layer of gram-negative organisms (2--7 nm, 1--3 layers) allows the CV-I complex to be extracted through the outer membrane disrupted by the organic solvent. (4) Counterstain: Safranin (red/pink) or basic fuchsin is applied for 30--60 seconds to stain the now colorless gram-negative organisms. Gram-positive organisms remain purple; gram-negative organisms appear red/pink. In addition to categorizing bacteria, the Gram stain evaluates the host inflammatory response: presence and quantity of polymorphonuclear leukocytes (PMNs/WBCs), squamous epithelial cells (SECs, indicating oropharyngeal contamination in respiratory specimens), mononuclear cells (CSF), and the relationship between organisms and cells (intracellular versus extracellular) \cite{murray2019manual}. The semi-quantitative assessment of organisms is reported as rare, few, moderate, or many per oil-immersion field (1,000x). The morphology (cocci in chains, clusters, pairs; diplococci; rods/bacilli; coccobacilli; branching filaments; spore formers; yeast forms) provides additional taxonomic information.

\section*{History}
The Gram stain was developed in 1884 by the Danish bacteriologist Hans Christian Gram (1853--1938). Gram was attempting to develop a method to visualize bacteria in tissue sections of lung from patients who had died of pneumonia. He discovered that certain bacteria retained the crystal violet-iodine stain after decolorization with alcohol, while others did not. The discovery was published in 1884 in a paper entitled "Über die isolierte Färbung der Schizomyceten in Schnitt- und Trockenpräparaten" (On the isolated staining of Schizomycetes in section and dried preparations) \cite{popescu1996gram}. The method was quickly adopted by Carl Friedländer and other bacteriologists and remains essentially unchanged after 140 years. The practical value of the Gram stain for guiding empiric antibiotic therapy is immense: broad-spectrum antibiotics can be narrowed based on the Gram stain result, and in critically ill patients (meningitis, septic shock), the Gram stain may be the only rapid diagnostic result available to guide initial antibiotic selection. In the 1970s--1980s, the systematic use of Gram stain for sputum specimen quality assessment (Bartlett's criteria: <10 SECs/LPF, >25 WBCs/LPF) was promoted and has since become universal practice in clinical microbiology \cite{bartlett1987diagnosis}.

\section*{Why This Test Is Ordered}
\begin{itemize}
  \item Rapid, presumptive identification of bacteria in clinical specimens from virtually any body site, guiding initial empiric antibiotic therapy while culture results are pending. Gram-positive and gram-negative organisms are susceptible to different classes of antibiotics; narrowing the differential at the Gram stain stage has significant antibiotic stewardship value
  \item STAT Gram stain of CSF in suspected bacterial meningitis: Gram-positive diplococci (\textit{S. pneumoniae}), gram-negative diplococci (\textit{N. meningitidis}), gram-positive rods (\textit{Listeria monocytogenes}), gram-negative rods (neonatal \textit{E. coli} or \textit{H. influenzae}---pleomorphic coccobacilli). The Gram stain result at 30--60 minutes after LP guides empiric antibiotic selection: vancomycin + ceftriaxone for gram-positive diplococci; ceftriaxone (or cefotaxime) for gram-negative diplococci; ampicillin + gentamicin or ceftriaxone for \textit{Listeria} in at-risk populations
  \item Assessment of sputum specimen adequacy: The presence of >10 squamous epithelial cells per low-power field (100x) indicates oropharyngeal contamination. The presence of >25 WBCs per LPF with <10 SECs indicates an adequate specimen from the lower respiratory tract
  \item Rapid diagnosis of gonococcal urethritis in symptomatic males: Gram stain of urethral discharge showing gram-negative intracellular diplococci within PMNs has sensitivity >95\% and specificity >99\%, allowing immediate diagnosis and treatment
  \item Evaluation of wound infections: The Gram stain can rapidly identify group A streptococcus (gram-positive cocci in chains), \textit{S. aureus} (gram-positive cocci in clusters), polymicrobial infection with mixed gram-positive and gram-negative organisms, or \textit{Clostridium perfringens} (gram-positive boxcar-shaped rods with few to no WBCs---suggestive of clostridial myonecrosis/gas gangrene)
  \item Evaluation of synovial fluid for septic arthritis: The presence of organisms on Gram stain (combined with WBC count typically >50,000 with >90\% neutrophils) confirms septic arthritis. A negative Gram stain does not exclude it (sensitivity 50--70\% for gonococcal arthritis, higher for staphylococcal and streptococcal septic arthritis)
  \item Rapid diagnosis of bacterial vaginosis: Gram stain with Nugent scoring evaluates the shift in vaginal flora from \textit{Lactobacillus}-predominant to BV-associated anaerobes
\end{itemize}

\section*{Primary vs Secondary Test}
Gram stain is a primary rapid test performed on virtually all clinical microbiology specimens. It is ordered reflexively when a culture is ordered for any specimen type. The Gram stain result is available within 1--2 hours of specimen receipt and provides preliminary guidance, while culture provides definitive identification and susceptibility.

\section*{How This Test Is Performed}
A thin smear of the specimen is prepared on a glass slide. For liquid specimens (CSF, body fluids, urine, BAL), a drop is placed on the slide and spread. For swabs, the swab is rolled gently onto the slide to avoid destroying cellular architecture. For thick or purulent specimens, the specimen is spread thinly. The slide is air-dried and then heat-fixed by passing through a flame (for most bacteria) or fixed in methanol (for CSF or when Gram staining for intracellular organisms) \cite{thomson2019specimen}. The four-step staining procedure described in Principle is performed. The slide is examined under oil immersion (1,000x magnification) by scanning multiple fields. The report includes: (1) Quantitation and type of WBCs (PMNs, mononuclear cells); (2) Quantitation of SECs (for respiratory specimens); (3) Presence, quantity, and Gram reaction/morphology of any organisms; (4) Intracellular or extracellular location of organisms. Quality control slides (gram-positive \textit{S. aureus} and gram-negative \textit{E. coli}) should be stained in parallel to ensure proper staining.

\section*{What Preparation Is Needed}
No patient preparation. The specimen collection (blood, CSF, swab, sterile fluid) is described in the respective chapters. For sputum Gram stain, the patient should rinse the mouth with water before expectoration. Slides must be fixed and stained within 2 hours of smear preparation to avoid bacterial overgrowth or autolysis.

\section*{Contraindications and Precautions}
\begin{itemize}
  \item The sensitivity of Gram stain for detecting bacteria is approximately 10$^{4}$--10$^{5}$ CFU/mL (organisms per mL of specimen). At lower concentrations, organisms will not be visualized, and a negative Gram stain does NOT exclude infection. CSF Gram stain in bacterial meningitis has a sensitivity of 60--90\% (highest for \textit{S. pneumoniae} and \textit{N. meningitidis}; lower for \textit{L. monocytogenes} and \textit{H. influenzae}). The Gram stain should NEVER be used alone to exclude bacterial meningitis---CSF cell count, glucose, protein, and culture (and increasingly, multiplex PCR meningitis/encephalitis panel) are essential. Gram stain misinterpretation is possible:
  \item (1) Over-decolorization: Gram-positive organisms may appear gram-negative (false gram-negative), especially if the ethanol step is prolonged.
  \item (2) Under-decolorization: Gram-negative organisms may retain crystal violet and appear gram-positive (false gram-positive).
  \item (3) Dead or damaged bacteria may stain variably or poorly.
  \item (4) Certain organisms have characteristic Gram stain variability: \textit{Acinetobacter} species may appear as gram-positive cocci in clinical specimens (they are actually gram-negative coccobacilli that resist decolorization). \textit{Clostridium} species may appear gram-negative in old cultures.
  \item (5) Small numbers of organisms may be missed due to sampling error or obscured by purulent material. The Gram stain cannot determine the species of bacteria (with rare exceptions like \textit{N. gonorrhoeae} intracellular diplococci, or \textit{Actinomyces} branching gram-positive filaments with sulfur granules). Culture or molecular identification is required.
\end{itemize}
\section*{Risks}
There are no risks from the Gram stain test itself (the stain is performed in the laboratory on the collected specimen). Risks are associated with specimen collection. The major risk is failure to detect organisms on Gram stain, leading to inappropriate delay in antibiotic therapy or inappropriate narrow-spectrum antibiotic selection.

\section*{What Happens After the Test}
Gram stain results are available STAT, typically within 1--2 hours of specimen receipt in the laboratory. The result is immediately communicated to the ordering clinician as a critical value when organisms are seen in a normally sterile site (CSF, blood, synovial fluid, pleural fluid, peritoneal fluid). Empiric antibiotic therapy is modified based on the Gram stain result: gram-positive cocci in clusters suggests staphylococci (add vancomycin for possible MRSA); gram-positive cocci in chains suggests streptococci (penicillin or ceftriaxone); gram-negative rods suggest Enterobacterales or \textit{Pseudomonas} (broad-spectrum beta-lactam: piperacillin-tazobactam, cefepime, or meropenem). Definitive identification and susceptibility follow at 48--72 hours.

\section*{Understanding Results}

\subsection*{Below Normal}
\begin{itemize}
  \item \textbf{No organisms seen:} Bacteria were not visualized on the Gram stain. This may indicate: absence of bacterial infection (viral, fungal, non-infectious process); bacterial infection with organism concentration below the limit of detection (<10$^{4}$ CFU/mL); effectively treated infection; or technical error (inadequate specimen, thick smear obscuring organisms, over-decolorization). This result does NOT exclude bacterial infection, and empiric antibiotics should be continued if clinical suspicion is high pending culture results
  \item \textbf{No WBCs seen:} Absence of an inflammatory response. May indicate: non-infectious process; infection at a very early stage; or severe immunosuppression (neutropenia) where PMNs are absent
  \item \textbf{Sputum rejected (>10 SECs/LPF):} Specimen is contaminated with oropharyngeal flora and is not representative of the lower respiratory tract. Repeat collection is recommended; culture may be performed but results should be interpreted with caution
\end{itemize}

\subsection*{Above Normal}
\begin{itemize}
  \item \textbf{Gram-positive cocci in clusters:} Characteristic of staphylococci (\textit{S. aureus}, coagulase-negative staphylococci). Empiric therapy: vancomycin (for possible MRSA in hospitalized or at-risk patients). MSSA can be de-escalated to nafcillin/cefazolin if susceptible
  \item \textbf{Gram-positive cocci in chains (or pairs):} Characteristic of streptococci (\textit{S. pneumoniae}: lancet-shaped diplococci; \textit{S. pyogenes}: chains; \textit{Enterococcus}: pairs/short chains). Empiric therapy: penicillin, ampicillin, or ceftriaxone (for pneumococcus); ampicillin or vancomycin (for enterococcus). Vancomycin should be added for suspected meningitis pending culture
  \item \textbf{Gram-negative diplococci (intracellular):} Characteristic of \textit{Neisseria} species (\textit{N. meningitidis} in CSF/blood; \textit{N. gonorrhoeae} in urethral/cervical exudate). In a symptomatic male, gram-negative intracellular diplococci in PMNs from urethral discharge is diagnostic of gonorrhea
  \item \textbf{Gram-negative rods:} Enterobacterales (\textit{E. coli}, \textit{Klebsiella}, \textit{Proteus}) or non-fermenters (\textit{Pseudomonas aeruginosa}, \textit{Acinetobacter}). Empiric therapy: broad-spectrum beta-lactam (cefepime, piperacillin-tazobactam, meropenem) based on local antibiogram and patient risk factors for resistance
  \item \textbf{Gram-positive rods:} \textit{Bacillus} species (large, boxcar-shaped, may have spores), \textit{Corynebacterium} (club-shaped, palisading/chinese-letter arrangement), \textit{Listeria monocytogenes} (small, regular rods; may be mistaken for \textit{S. pneumoniae} on Gram stain---the "diphtheroid mimic"). In CSF, gram-positive rods suggest \textit{Listeria} (ampicillin added to empiric regimen)
  \item \textbf{Gram-variable or mixed morphology:} May suggest polymicrobial infection, contamination, or organisms with unusual Gram-staining characteristics (e.g., \textit{Acinetobacter}, \textit{Clostridium} in older cultures)
\end{itemize}

\section*{Normal Results}
\textbf{No organisms seen.} \textbf{Polymorphonuclear leukocytes (PMNs):} appropriate number for the specimen site. For sputum: \textbf{Squamous epithelial cells <10 per LPF; WBC >25 per LPF} (adequate specimen). For CSF: \textbf{WBC 0--5 cells/$\mu$L.} Sterile body fluids: \textbf{No organisms seen; no WBCs or few WBCs.}

\section*{Follow-up Tests}
\begin{itemize}
  \item \textbf{Bacterial culture with antimicrobial susceptibility testing:} Required for all specimens submitted for Gram stain; the Gram stain provides preliminary information, but culture and AST are needed for definitive species identification, confirmation of the preliminary Gram stain findings, and antibiotic susceptibility
  \item \textbf{MALDI-TOF MS of culture isolate:} For rapid species-level identification from growth on the culture plates (minutes from colony)
  \item \textbf{PCR or multiplex panel (e.g., BioFire FilmArray Meningitis/Encephalitis Panel, Respiratory Panel):} When Gram stain is negative but clinical suspicion is high, or when the Gram stain is positive but the organism is fastidious and may not grow in culture (e.g., partially treated meningitis; \textit{Kingella kingae} in pediatric septic arthritis). PCR may detect nonviable organisms after antibiotics
  \item \textbf{Acid-fast stain (Ziehl-Neelsen, auramine-rhodamine):} When organisms are not visualized on Gram stain and mycobacterial or nocardial infection is suspected. \textit{Nocardia} and \textit{Mycobacterium} are gram-positive but are poorly visualized on Gram stain; they are better detected by AFB stain (modified Kinyoun for \textit{Nocardia})
\end{itemize}

\section*{References}
\bibliographystyle{unsrtnat}
\bibliography{references}

\clearpage
\section*{Regulatory Status and Authorization Date}
This chapter describes a test category, not a specific commercial product. FDA, CDSCO, EU IVDR, and MHRA approval or registration dates therefore cannot be assigned without the assay manufacturer, product name, instrument, and intended use. For laboratory-developed tests, an individual agency approval date may not exist. See the Preface for the interpretation of regulatory status.

\clearpage
